\documentclass[unicode,11pt,a4paper,oneside,numbers=endperiod,openany]{scrartcl}

\usepackage{minted}
\usepackage{amsmath}

\usepackage{ifthen}
\usepackage[utf8]{inputenc}
\usepackage{graphics}
\usepackage{graphicx}
\usepackage{hyperref}

\pagestyle{plain}
\voffset -5mm
\oddsidemargin  0mm
\evensidemargin -11mm
\marginparwidth 2cm
\marginparsep 0pt
\topmargin 0mm
\headheight 0pt
\headsep 0pt
\topskip 0pt        
\textheight 255mm
\textwidth 165mm

\newcommand{\duedate} {}
\newcommand{\setduedate}[1]{%
\renewcommand\duedate {Due date:~ #1}}
\newcommand\isassignment {false}
\newcommand{\setassignment}{\renewcommand\isassignment {true}}
\newcommand{\ifassignment}[1]{\ifthenelse{\boolean{\isassignment}}{#1}{}}
\newcommand{\ifnotassignment}[1]{\ifthenelse{\boolean{\isassignment}}{}{#1}}

\newcommand{\assignmentpolicy}{
\begin{table}[h]
\begin{center}
\scalebox{0.8} {%
\begin{tabular}{|p{0.02cm}p{16cm}|}
\hline
&\\
\multicolumn{2}{|c|}{\Large\textbf{HPC Lab for CSE 2024 ---  Submission Instructions}}\\
\multicolumn{2}{|c|}{\large\textbf{(Please, notice that following instructions are mandatory: }}\\
\multicolumn{2}{|c|}{\large\textbf{submissions that don't comply with, won't be considered)}}\\
&\\
\textbullet & Assignments must be submitted to \href{https://moodle-app2.let.ethz.ch/course/view.php?id=22516}{Moodle} (i.e. in electronic format).\\
\textbullet & Provide both executable package and sources (e.g. C/C++ files, Matlab). 
If you are using libraries, please add them in the file. Sources must be organized in directories called:\\
\multicolumn{2}{|c|}{\textit{Project\_number\_lastname\_firstname}}\\
& and  the  file must be called:\\
\multicolumn{2}{|c|}{\textit{project\_number\_lastname\_firstname.zip}}\\
\multicolumn{2}{|c|}{\textit{project\_number\_lastname\_firstname.pdf}}\\
\textbullet &  The TAs will grade your project by reviewing your project write-up, and looking at the implementation 
                 you attempted, and benchmarking your code's performance.\\

\textbullet & You are allowed to discuss all questions with anyone you like; however: (i) your submission must list anyone you discussed problems with and (ii) you must write up your submission independently.\\
\hline
\end{tabular}
}
\end{center}
\end{table}
}
\newcommand{\punkte}[1]{\hspace{1ex}\emph{\mdseries\hfill(#1~\ifcase#1{Points}\or{Points}\else{Points}\fi)}}


\newcommand\serieheader[6]{
\thispagestyle{empty}%
\begin{flushleft}
\includegraphics[width=0.4\textwidth]{ETHlogo_13}
\end{flushleft}
  \noindent%
  {\large\ignorespaces{\textbf{#1}}\hspace{\fill}\ignorespaces{ \textbf{#2}}}\\ \\%
  {\large\ignorespaces #3 \hspace{\fill}\ignorespaces #4}\\
  \noindent%
  \bigskip
  \hrule\par\bigskip\noindent%
  \bigskip {\ignorespaces {\Large{\textbf{#5}}}
  \hspace{\fill}\ignorespaces \large \ifthenelse{\boolean{\isassignment}}{\duedate}{#6}}
  \hrule\par\bigskip\noindent%  \linebreak
 }

\makeatletter
\def\enumerateMod{\ifnum \@enumdepth >3 \@toodeep\else
      \advance\@enumdepth \@ne
      \edef\@enumctr{enum\romannumeral\the\@enumdepth}\list
      {\csname label\@enumctr\endcsname}{\usecounter
        {\@enumctr}%%%? the following differs from "enumerate"
	\topsep0pt%
	\partopsep0pt%
	\itemsep0pt%
	\def\makelabel##1{\hss\llap{##1}}}\fi}
\let\endenumerateMod =\endlist
\makeatother




\usepackage{textcomp}






\begin{document}


\setassignment
\setduedate{11 March 2024, 23:59}

\serieheader{High-Performance Computing Lab for CSE}{2024}
{Student: Benedict Armstrong}
{Discussed with: FULL NAME}{Solution for Project 1a}{}
\newline



% \assignmentpolicy

\section{Euler warm-up [10 points]}

\subsection{Module System}
% On the cluster, we provide many centrally installed software and for some software even multiple versions. To configure the environment for a particular software version, we use modules. Modules configure your current computing environment (PATH, LD_LIBRARY_PATH, MANPATH, etc.) to make sure all required binaries and libraries are found.

The module system allows Euler users to quickly and easily configure their environment to use centrally installed software package. A detailed description can be found in the \href{https://scicomp.ethz.ch/wiki/Modules_and_applications}{Module System} documentation.

There are two systems currently in use. The older system is called \texttt{Environment Modules} and the newer system is called \texttt{LMOD Modules}. All new software installations are done with LMOD Modules.

% Code examples
\begin{listing}[!ht]
    \begin{minted}{bash}
    # List all available modules
    module avail
    
    # Load a module
    module load <module_name>
    
    # list all loaded modules
    module list
    \end{minted}
    \caption{Module System}
    \label{lst:module_system}
\end{listing}

\subsection{SLURM}

The Euler cluster uses SLURM to manage and schedule jobs. To run a job on the cluster, you need to submit a job script to the SLURM scheduler. A detailed description can be found in the \href{https://scicomp.ethz.ch/wiki/Job_management_with_SLURM}{SLURM} documentation.

\subsection{Hello Euler!}

We start by compiling and running a simple C program on the Euler cluster. The program is called \texttt{hello\_euler.cpp} and should print "\texttt{Host name: <hostname>}" to standard out.

To run the compiled program on the cluster, we need to submit a job script to the SLURM scheduler. The job script is called \texttt{hello\_euler.slurm} and should look like this:

\begin{listing}[!ht]
    \inputminted{bash}{../01a/hello_euler/hello_euler_1.sh}
    \caption{Job script for running hello\_euler.cpp}
    \label{lst:hello_euler}
\end{listing}

The job can then be submitted to the SLURM scheduler with the following command:

\begin{listing}[!ht]
    \begin{minted}{bash}
    sbatch hello_euler.sh
    \end{minted}
    \caption{Submitting a job to the SLURM scheduler}
    \label{lst:sbatch}
\end{listing}

The code and output can be found in the \texttt{hello\_euler} directory.

\subsection{Multiple Nodes}

We can run the same code on multiple nodes using the following job script:


\begin{listing}[!ht]
    \inputminted{bash}{../01a/hello_euler/hello_euler_2.sh}
    \caption{Job script for running hello\_euler.cpp on multiple nodes}
    \label{lst:hello_euler_2}
\end{listing}

Where we set the number of nodes to 2 and the number of tasks to 2. The output can be found in the \texttt{hello\_euler\_2.out} file.

\section{Performance characteristics [50 points]}

\subsection{Peak performance}

The peak performance of a cluster can be calculated using the following formula:


\begin{align*}
    p_{core}    & = n_{super} \times n_{FMA} \times n_{SIMD} \times f_{core} \\
    p_{CPU}     & = n_{core} \times p_{core}                                 \\
    p_{nodes}   & = n_{nodes} \times p_{CPU}                                 \\
    p_{cluster} & = p_{nodes} \times n_{cluster}
\end{align*}


in our case this is equal to:

\begin{table}
    \centering

    \begin{tabular}{|c|c|c|c}
        \hline
        \textbf{Parameter} & \textbf{EPYC\_7H12} & \textbf{EPYC\_7763} \\
        \hline
        $n_{super}$        & 2                   & 2                   \\
        \hline
        $n_{FMA}$          & 2                   & 2                   \\
        \hline
        $n_{SIMD}$         & 256                 & 256                 \\
        \hline
        $f_{core}$         & 2.25GHz             & 2.45GHz             \\
        \hline
        $n_{core}$         & 64                  & 64                  \\
        \hline
        $n_{nodes}$        & 2                   & 2                   \\
        \hline
        $n_{cluster}$      & 1                   & 1                   \\
        \hline
        $p_{core}$         & 2.25GHz             & 2.45GHz             \\
        \hline
        $p_{CPU}$          & 144GHz              & 156.8GHz            \\
        \hline
        $p_{nodes}$        & 288GHz              & 313.6GHz            \\
        \hline
        $p_{cluster}$      & 288GHz              & 313.6GHz            \\
        \hline
    \end{tabular}
    \caption{Peak performance of the EPYC\_7H12 and EPYC\_7763 nodes}
    \label{tab:peak_performance}
\end{table}

\subsection{Memory Hierarchies}

The output of running \texttt{lscpu} and \texttt{hwloc-ls} can be found in the \texttt{memory\_hierarchies} directory. As in the example in the assignment there are also two PDFs detailing the memory hierarchy of the EPYC\_7H12 and EPYC\_7763 nodes. In summary both nodes have 8 NUMA nodes, with 8 cores per NUMA node. More information on NUMA can easily be found in the \href{https://en.wikipedia.org/wiki/Non-uniform_memory_access}{Wikipedia} page. Basically it means that the nodes have faster access to their specific part of the shared memory. The rest of the numbers can easily be read out of the two PDFs detailing the memory hierarchy.

\subsubsection{Cache and main memory size}

% Tables for cache (L1-L3) and main memory size for both nodes (EPYC_7H12 and EPYC_7763) and if they are shared or not

\begin{table}[h]
    \centering
    \begin{tabular}{|c|c|c|}
        \hline

        \textbf{Cache} & \textbf{EPYC\_7H12} & \textbf{EPYC\_7763} \\
        \hline
        L1d            & 32KB                & 32KB                \\
        \hline
        L1i            & 32KB                & 32KB                \\
        \hline
        L2             & 512KB               & 512KB               \\
        \hline
        L3             & 16MB                & 32MB                \\
        \hline
        NUMA           & 63GB                & 31GB                \\
        \hline
        Total Machine  & 502GB               & 248GB               \\
        \hline
    \end{tabular}
    \caption{Cache and main memory size for both nodes}
    \label{tab:cache_main_memory}
\end{table}

\subsection{Bandwidth: STREAM benchmark}

\subsection{Performance model: A simple roofline model}


\end{document}
